\documentclass[11pt]{article}

\usepackage[margin=1in]{geometry}
\usepackage{amsmath}
\usepackage{amssymb}
\usepackage{booktabs}

\title{A Modular Base-$e$ Parameterization of Lepton and Hadron Mass Ratios}
\author{I. Mikhailov}
\date{May 2026}

\begin{document}
\maketitle

\begin{abstract}
We present a compact phenomenological parameterization of selected lepton and
hadron mass ratios using the constants $e$, $\pi$, the inverse fine-structure
constant $\alpha^{-1}$, the topological index $N=5$, the bosonic-string
reference dimension $D=26$, and the Fibonacci number $F_{26}=121393$.  The
construction is not derived from a Lagrangian and is not proposed as a
replacement for the Standard Model.  Its purpose is narrower: to record a
low-complexity set of algebraic formulae, their numerical residuals, and the
places where the ansatz fails or requires a non-standard interpretation.
\end{abstract}

\section{Discrete Basis}

The working basis is
\begin{equation}
    P=\{2,3,5,13,137\}, \qquad N=5,\qquad D=26,\qquad F_{26}=121393.
\end{equation}
The integer $137$ is treated as a discrete anchor.  In the exploratory
base-$e$ notation used in this project,
\begin{equation}
    137_e = 21100.0210102\ldots_e ,
\end{equation}
with the first fractional digit pattern
\begin{equation}
    d_{-1},d_{-2},d_{-3},d_{-4},d_{-5},d_{-6},d_{-7}
    =
    0,2,1,0,1,0,2.
\end{equation}
The earlier digit moment
\begin{equation}
    I_5=\sum_{k=1}^{5}k^2d_{-k}=42
\end{equation}
has the equivalent topological form
\begin{equation}
    I_5=2(D-N)=2(26-5)=42.
\end{equation}
The core compression operator is
\begin{equation}
    q = \frac{N(N-2)}{e^4\pi^3}
      = \frac{15}{e^4\pi^3}
      =0.008860611874290873.
\end{equation}
We also use
\begin{equation}
    s=\frac{\alpha}{2\pi}, \qquad
    \delta\phi=\cos(1/N)-\cos(2/N)=0.05900558383835652.
\end{equation}

\section{Charged Lepton Sector}

The electron is used as the normalization anchor.  The project-level
normalization is
\begin{equation}
    C_e=\frac{1}{1-s\left(N+\frac{2}{eN}\right)},
\end{equation}
and the electron scale is written as
\begin{equation}
    m_e = 10e^{11}\cos^2(2/N)C_e .
\end{equation}
This equation is a normalization convention in the present note.

The muon ratio is
\begin{equation}
    \frac{m_\mu}{m_e}
    =
    \alpha^{-1}\left(\frac{3}{2}+q\right).
\end{equation}
The tau ratio is
\begin{equation}
    \frac{m_\tau}{m_e}
    =
    \alpha^{-1}\left[N^2+qf_\tau\right],
\end{equation}
where
\begin{equation}
    f_\tau = 2(D-N)+2e^{-2}+2e^{-7}.
\end{equation}

\section{Hadron and Baryon Benchmarks}

For the neutral pion, the baseline expression is
\begin{equation}
    \frac{m_{\pi^0}}{m_e}
    =
    \alpha^{-1}
    \left[
        2-\delta\phi
        -
        q\left(1+\frac{1}{N}+\frac{1}{\pi}\right)
    \right].
\end{equation}
The Phase 9 audit showed that the literal replacement $q\rightarrow117/137$
fails by changing the scale of the expression.  The retained compact
modulation uses the invariant
\begin{equation}
    117=DN-13=9\cdot13
\end{equation}
only as a small correction:
\begin{equation}
    q_{117}=q\left(1-\frac{2^2}{117\cdot137}\right).
\end{equation}
The accepted pion benchmark is therefore
\begin{equation}
    \frac{m_{\pi^0}}{m_e}
    =
    \alpha^{-1}
    \left[
        2-\delta\phi
        -
        q_{117}\left(1+\frac{1}{N}+\frac{1}{\pi}\right)
    \right].
\end{equation}

For the proton ratio, the low-complexity Phase 9 expression is
\begin{equation}
    K_p=\frac{M_p}{m_e}
    =
    \frac{2F_{26}}{DN}\cos(1/N)
    \left[
        1+es+\frac{6}{137^2(D+N+2^2)}
    \right].
\end{equation}
The integer $6$ is the Euler-phi gap
\begin{equation}
    \phi(137)-DN=136-130=6.
\end{equation}
This should be read as a compact empirical candidate, not as a derivation from
QCD.

\section{Predicted Sterile Benchmarks}

The charged heavy-lepton extrapolation is
\begin{equation}
    \frac{m_{\tau'}}{m_e}
    =
    \alpha^{-1}\left[D^2+q(DN)\right],
\end{equation}
which gives
\begin{equation}
    m_{\tau'}=47.417730830364825~\mathrm{GeV}.
\end{equation}
This is above the project-level $45~\mathrm{GeV}$ reference threshold, but
below the conventional sequential charged-heavy-lepton limit of about
$100.8~\mathrm{GeV}$.  It therefore cannot be interpreted as an ordinary
Standard-Model-like fourth charged lepton without additional assumptions.

The sterile neutral benchmark is
\begin{equation}
    m_{\nu_{\tau'}}
    =
    m_e
    \left(\frac{F_{26}}{\alpha^{-1}}\right)
    \left[
        1+\frac{DN}{\pi}\delta\phi(1-s)
    \right],
\end{equation}
which gives
\begin{equation}
    m_{\nu_{\tau'}}=1.5566463427697075~\mathrm{GeV}.
\end{equation}
This mass is kinematically cold, but the relic abundance, lifetime, production
history, and visible decay channels are not determined by the mass formula.
The benchmark is conditionally viable only for a fully sterile or otherwise
sufficiently decoupled state.  A simple Fermi-LAT/Galactic-Center audit did not
support a direct gamma-ray-excess interpretation: the phase-modulated
annihilation estimates are far below the nominal
$\sim2\times10^{-26}~\mathrm{cm^3\,s^{-1}}$ scale.

\section{Numerical Summary}

\begin{table}[h]
\centering
\caption{Mass-ratio benchmarks recorded by the model.  Residuals are quoted
relative to the reference values used in the project audit.}
\begin{tabular}{llll}
\toprule
Observable & Formula family & Model value & Residual \\
\midrule
$m_\mu/m_e$ & $\alpha^{-1}(3/2+q)$ & $206.76822$ & $-0.29~\mathrm{ppm}$ \\
$m_\tau/m_e$ & $\alpha^{-1}(N^2+qf_\tau)$ & $3477.22820$ & $-0.022~\mathrm{ppm}$ \\
$m_{\pi^0}$ & $q_{117}$ pion branch & $134.976797~\mathrm{MeV}$ & $-0.020995~\mathrm{ppm}$ \\
$M_p/m_e$ & $\phi(137)$ gap branch & $1836.152490$ & $-0.098658~\mathrm{ppm}$ \\
$m_{\tau'}$ & $\alpha^{-1}[D^2+qDN]$ & $47.417731~\mathrm{GeV}$ & benchmark \\
$m_{\nu_{\tau'}}$ & sterile branch & $1.556646~\mathrm{GeV}$ & benchmark \\
\bottomrule
\end{tabular}
\end{table}

\section{Conclusion}

The compact relations collected here are best viewed as a modular
parameterization of selected particle scales.  The formulas are numerically
tight and reproducible in the accompanying code, but they remain
phenomenological until connected to a dynamical field theory, collider
constraints, and standard hadronic calculations.  The useful output of the
exercise is therefore not a claim of replacement physics, but a falsifiable
matrix of low-complexity numerical benchmarks.

\end{document}
